\chapter{Estrutura do reposit\'orio de c\'odigo}
\label{ap:repo}

O reposit\'orio Git que acompanha este trabalho est\'a organizado em
cinco diret\'orios de primeiro n\'ivel, cada um cobrindo uma fase
distinta do ciclo experimental: \texttt{Training/} concentra os
\emph{scripts} de constru\c{c}\~ao de grafos, treinamento e
avalia\c{c}\~ao; \texttt{Execution/} arquiva os artefatos persistidos
(pesos, resultados, figuras); \texttt{Material/} cont\'em os
\emph{datasets} brutos e os arquivos do TCC; \texttt{Interface/}
abriga a aplica\c{c}\~ao web demonstrativa do
Ap\^endice~\ref{ap:aplicacao}; e o n\'ivel raiz cont\'em os
documentos de controle do trabalho.

% ============================================================
\section{\texttt{Training/} --- \emph{scripts} experimentais}
\label{sec:ap_training}

Subdivide-se em tr\^es fases hist\'oricas, das quais a relevante para
a vers\~ao final \'e \path{Training/03_Mega_Research/}. Os
diret\'orios \path{01_BlueSky_Pipe/} e \path{02_UPFD_Benchmark/}
correspondem a itera\c{c}\~oes anteriores do \emph{pipeline} e s\~ao
mantidos para refer\^encia hist\'orica; suas funcionalidades est\~ao
consolidadas, em forma refatorada, em
\path{03_Mega_Research/}.

A pasta \path{Training/03_Mega_Research/} cont\'em $30$
\emph{scripts} numerados de \texttt{00} a \texttt{32} (com sufixos
\texttt{27b}, \texttt{27c}, \texttt{31}), organizados pelo papel no
\emph{pipeline}:

\begin{sloppypar}
\begin{description}\setlength{\itemsep}{2pt}
\item[Constru\c{c}\~ao de grafos:] \path{00_construir_grafos_fakenewsnet.py},
\path{01_baixar_dataset_hf.py}, \path{02_construir_grafos_bluesky.py}.

\item[Treinamento e benchmark de arquiteturas:]
\path{03_treinar_gat.py}, \path{04_comparar_gcn_gat.py},
\path{05_treinar_sage.py}, \path{06_comparar_gcn_gat_sage.py},
\path{07_upfd_benchmark_triplo.py},
\path{13_benchmark_upfd_oficial.py}.

\item[Avalia\c{c}\~ao estat\'istica:] \path{09_teste_significancia.py}
(\texttt{ttest\_rel} pareado por \emph{fold}); \path{15_analise_estrutural.py}
(Cohen's $d$, T5/F17).

\item[\emph{Baselines} e diagn\'osticos:] \path{10_baseline_textual.py}
(LogReg-BERT e RF, T1), \path{11_diagnostico_confound.py}
(RF sobre \texttt{num\_nodes}), \path{20_textual_vs_topologico.py}
(concord\^ancia, T6/F3).

\item[\emph{Ablation} de \emph{features}:]
\path{12_ablation_intra_encoding.py},
\path{14_topologia_sem_texto.py} (T4/F18, achado central).

\item[Explicabilidade:] \path{16_gnn_explainer_upfd.py} (F19, F12).

\item[Persist\^encia de modelos:] \path{17_persistir_modelos_finais.py}
(salva LogReg-BERT, RF estrutural e SAGE estrutural em
\path{Execution/weights/}).

\item[Bluesky:] \path{18_analise_bluesky_crossfeed.py},
\path{19_aplicar_modelo_bluesky.py},
\path{22_concordancia_bluesky.py},
\path{23_visualizar_threads_bluesky.py},
\path{26_rq3_multilingual.py} (T11/F14).

\item[Consolida\c{c}\~ao e publica\c{c}\~ao:]
\path{21_bloco_vs_mini.py} (T7/F4),
\path{24_consolidar_tcc.py} (gera o
\path{INDICE.md}), \path{25_publicar_huggingface.py},
\path{27_figuras_comparativas.py},
\path{27b_repintar_figuras_orientador.py} (paleta sem\^antica),
\path{27c_matrizes_confusao_mono.py} (F33--F35).

\item[\emph{Fase 9}, estratifica\c{c}\~ao do erro:]
\path{28_distribuicao_e_estratificacao.py} (F25--F27),
\path{29_outliers.py} (T12/F29/F30),
\path{30_cascatas_profundas.py} (T13/F28/F31/F32).

\item[Anota\c{c}\~ao de figuras:] \path{31_anotar_figuras_cap4.py}
(sobrep\~oe c\'irculos+texto em PNGs para vers\~oes destacadas em
\path{Imagens/anotadas/}).
\end{description}
\end{sloppypar}

Os \emph{scripts} de defini\c{c}\~ao de arquitetura
(\path{gcn_model.py}, \path{gat_model.py}, \path{sage_model.py})
e o utilit\'ario \path{gerar_folds.py} ficam na raiz de
\path{03_Mega_Research/}.

% ============================================================
\section{\texttt{Execution/} --- artefatos persistidos}
\label{sec:ap_execution}

\begin{description}\setlength{\itemsep}{2pt}
\item[\texttt{Execution/weights/}] Pesos persistidos dos tr\^es
classificadores finais empregados no Ap\^endice~\ref{ap:aplicacao}
(LogReg-BERT, RF estrutural e SAGE estrutural), produzidos por
\path{17_persistir_modelos_finais.py}. Os tamanhos t\'ipicos
constam da Tabela~\ref{tab:hp_ambiente}.

\item[\texttt{Execution/results/}] Sa\'idas tabulares (CSVs) e arquivos
de relat\'orio (TXT) de cada experimento. Cada \emph{script} de
\path{03_Mega_Research/} grava aqui o conjunto de m\'etricas
brutas que alimenta as tabelas e figuras do trabalho. A
\path{figuras_tcc/} subdiret\'orio cont\'em c\'opias finais das
figuras incorporadas no TCC.

\item[\texttt{Execution/scripts/}] Utilit\'arios de orquestra\c{c}\~ao
(execu\c{c}\~ao em lote, limpeza de cache, sincroniza\c{c}\~ao com
Hugging Face).
\end{description}

% ============================================================
\section{\texttt{Material/} --- \emph{datasets} e o TCC}
\label{sec:ap_material}

\begin{description}\setlength{\itemsep}{2pt}
\item[\texttt{Material/politifact/}, \texttt{gossipcop/}] CSVs brutos
do FakeNewsNet baixados do reposit\'orio KaiDMML~\cite{shu2020fakenewsnet}.

\item[\texttt{Material/gossipcop\_pyg/}] Vers\~ao do GossipCop
constru\'ida em formato \emph{PyTorch Geometric}.

\item[\texttt{Material/LIAR/}] \emph{Dataset} alternativo usado em
ensaios iniciais; n\~ao integrado \`a vers\~ao final.

\item[\texttt{Material/GNN\_TCC\_atualizado/}] \textbf{Fonte canonica
deste TCC}: arquivos \texttt{.tex}, \texttt{.bib}, figuras e tabelas.
A compila\c{c}\~ao ocorre nesta pasta. O \path{abntex2-alf.bst}
e o \emph{stub} \texttt{abntex2.cls} acompanham o reposit\'orio para
garantir compila\c{c}\~ao em ambientes sem o pacote \texttt{abntex2}
instalado.

\item[\texttt{Material/GNN\_TCC\_extraido/}] Vers\~oes anteriores do
TCC, mantidas para refer\^encia hist\'orica.
\end{description}

% ============================================================
\section{\texttt{Interface/} --- aplica\c{c}\~ao web demonstrativa}
\label{sec:ap_interface}

\begin{description}\setlength{\itemsep}{2pt}
\item[\texttt{Interface/frontend/}] Aplica\c{c}\~ao Next.js +
FastAPI documentada no Ap\^endice~\ref{ap:aplicacao}. O \emph{backend}
exp\~oe o \emph{endpoint} \texttt{/api/result} que recebe um
\emph{post} candidato e devolve os tr\^es \emph{scores} (textual,
topol\'ogico, combinado) acompanhados de \emph{flag} de
concord\^ancia.

\item[\texttt{Interface/supabase/}] Configura\c{c}\~ao de persist\^encia
opcional para hist\'orico de consultas (n\~ao essencial \`a demonstra\c{c}\~ao).
\end{description}

% ============================================================
\section{Documentos de controle (raiz do reposit\'orio)}
\label{sec:ap_docs}

A raiz do reposit\'orio cont\'em documentos de controle do trabalho
que n\~ao integram o texto do TCC mas governam a escrita e a
auditoria:

\begin{description}\setlength{\itemsep}{2pt}
\item[\texttt{FATOS\_TCC.md}] Cat\'alogo de fatos verificados (n\'umeros,
m\'etricas, decis\~oes), com estado (\(\checkmark\) verificado,
\(\triangle\) discrepante, \(?\) a verificar) e lista de seis
discrep\^ancias entre o \emph{draft} antigo e os dados reais (D1--D6).

\item[\texttt{PENDENCIAS\_TCC.md}] Decis\~oes pendentes e mudan\c{c}as
planejadas mas n\~ao executadas; auditoria de cada decis\~ao
estrutural do trabalho.

\item[\texttt{REVISAO\_CARA\_DE\_IA.md}] Itens pendentes de revis\~ao
estil\'istica gerados por an\'alise adversarial autom\'atica.

\item[\texttt{HISTORIA\_FASE9.md}, \texttt{MUDANCAS\_FASE9.md}] Documentos
internos sobre a evolu\c{c}\~ao da an\'alise estratificada (Fase 9).

\item[\texttt{GUIA\_COAUTORES\_TCC.md}, \texttt{PROMPT\_IA\_TCC.md}]
Conven\c{c}\~oes de escrita para coautores e para sess\~oes de
assist\^encia automatizada.
\end{description}

% ============================================================
\section{Reprodu\c{c}\~ao integral do \emph{pipeline}}
\label{sec:ap_reproducao}

Sob o ambiente documentado no Ap\^endice~\ref{ap:hiperparametros}
(Python 3.12, PyTorch 2.10 + CPU, semente universal 42), a execu\c{c}\~ao
sequencial dos \emph{scripts} numerados em
\path{Training/03_Mega_Research/} regenera, na ordem, todos os
artefatos catalogados no Ap\^endice~\ref{ap:indice}. A reprodu\c{c}\~ao
n\~ao \'e \emph{bit-exact} entre m\'aquinas distintas devido ao
n\~ao-determinismo de \texttt{scatter}/\texttt{gather} do
\emph{PyTorch Geometric} (caveat documentado em \S\ref{sec:protocolo});
as conclus\~oes qualitativas s\~ao robustas a essa varia\c{c}\~ao.

A vers\~ao desta tese registrada no Git em \texttt{main} corresponde
ao \emph{commit} de submiss\~ao; futuras itera\c{c}\~oes do trabalho
ser\~ao publicadas em \emph{branches} marcadas no mesmo reposit\'orio.
